\section{Price Discovery After Scheduled Releases}\label{sec:discovery}

This section asks how much of the response to a release arrives after the first minute and
whether the surprise predicts it. We define, for a release at minute $T$ and horizon $h$, the
reaction from the last pre-release price and the drift from the first post-release price,
\begin{equation}\label{eq:decomp}
  R_{T,h} = \ln P_{T+h-1} - \ln P_{T-1}, \qquad
  D_{T,h} = \ln P_{T+h-1} - \ln P_{T},
\end{equation}
where $P_\tau$ is the close of the bar labelled $\tau$. The reaction is the sum of the first-minute
jump $R_{T,1}$ and the drift, $R_{T,h}=R_{T,1}+D_{T,h}$. Only the drift can be traded.

\subsection{The jump and the drift}

Figure~\ref{fig:cpi_paths} shows the average cumulative return around CPI releases from 2016 to
2026, separately for hot prints (headline surprise above $+0.5$ standard deviations, 38 releases)
and cool prints (below $-0.5$, 36 releases). Three facts stand out. First, all four markets jump in
the direction theory predicts within the first minute: after hot prints \ES\ falls about 37 bps,
\NQ\ about 55 bps, \YM\ about 31 bps and \ZN\ about 19 bps; after cool prints they rise by about 30,
42, 26 and 23 bps.
Second, the equity indices keep moving in the same direction afterwards. After hot prints \ES\
falls a further 15--18 bps, to about $-55$ bps 30 minutes after the release, and \NQ\ a further 20
bps, to about $-77$ bps; \YM\ falls a further 17 bps, to about $-48$ bps. After cool prints the
continuation is smaller, about 3--6 bps in the three equity indices. Third,
\ZN\ continues by only 4--5 bps after hot prints and slightly reverses after cool prints. Before the
release, all four markets are flat on average: there is no systematic pre-positioning in the 30
minutes before CPI.

\begin{figure}[htbp]
  \centering
  \includegraphics[width=\linewidth]{fig02_cpi_event_paths_4m}
  \caption{Average cumulative return around CPI releases, 2016--2026, relative to the last
  pre-release price, for hot prints (headline surprise above $+0.5$ SD) and cool prints (below
  $-0.5$ SD), \ES, \NQ, \YM\ and \ZN. Shaded: the window a strategy can trade (one to 30 minutes after the release).}
  \label{fig:cpi_paths}
\end{figure}

\subsection{Surprises predict the reaction and the drift}

For each release family we regress the reaction and the drift on the standardized surprise
components with heteroskedasticity-robust (HC3) standard errors. Table~\ref{tab:eventstudy}
reports the results for CPI, NFP and PCE. CPI is the release that matters in all four markets: its
headline surprise explains 27--30\% of the five-minute reaction, with the expected sign. A
one-standard-deviation hot surprise lowers \ES\ by 20.8 bps, \NQ\ by 27.7 bps and \ZN\ by 9.9 bps
over five minutes. \NQ\ reacts about 35\% more than \ES, consistent with the higher rate
sensitivity of long-duration technology stocks, and \YM\ about 8\% less (19.2 bps, $R^2$ 0.27), with
the same drift as \ES\ after the first minute ($-5.3$ against $-5.7$ bps). NFP moves prices almost as much as CPI on average
(30 bps in \ES), but in a direction the surprise hardly explains ($R^2$ of 1--3\%), because
payrolls, the unemployment rate and wages often point in different directions.

\begin{table}[htbp]
  \centering\small
  \caption{Response to CPI, NFP and PCE surprises, 2016--2026.}\label{tab:eventstudy}
  \begin{threeparttable}
  \begin{tabular}{llccccc}
    \toprule
    & & Mean $|R_{T,30}|$ & \multicolumn{2}{c}{5-minute reaction} & \multicolumn{2}{c}{Drift $T{+}1\to T{+}30$} \\
    \cmidrule(lr){4-5}\cmidrule(lr){6-7}
    Market & Release & (bps) & $\beta$ (bps/SD) & $R^2$ & $\beta$ (bps/SD) & $t$ \\
    \midrule
    \ES & CPI & 38.5 & $-20.8$ & 0.27 & $-5.7$  & $-1.4$ \\
    \NQ & CPI & 53.6 & $-27.7$ & 0.30 & $-4.6$  & $-0.9$ \\
    \ZN & CPI & 22.5 & $-9.9$  & 0.30 & $-0.06$ & $-0.04$ \\
    \YM & CPI & 33.2 & $-19.2$ & 0.27 & $-5.3$  & $-1.4$ \\
    \midrule
    \ES & NFP & 29.6 & $+2.8$ & 0.03 & $+0.2$ & 0.1 \\
    \NQ & NFP & 36.0 & $+2.5$ & 0.01 & $-0.6$ & $-0.4$ \\
    \ZN & NFP & 23.6 & $-1.5$ & 0.03 & $+0.1$ & 0.1 \\
    \YM & NFP & 27.9 & $+2.8$ & 0.07 & $+0.6$ & 0.4 \\
    \midrule
    \ES & PCE & 14.9 & $-3.8$ & 0.10 & $+0.4$ & 0.2 \\
    \NQ & PCE & 18.5 & $-5.2$ & 0.10 & $+1.3$ & 0.5 \\
    \ZN & PCE & 8.3  & $-0.9$ & 0.04 & $+0.5$ & 0.5 \\
    \YM & PCE & 11.9 & $-2.9$ & 0.09 & $-0.6$ & $-0.4$ \\
    \bottomrule
  \end{tabular}
  \begin{tablenotes}\footnotesize
    \item Notes: $\beta$ is the coefficient on the headline surprise in a regression of the
    return on all surprise components of the release (CPI: headline and core; NFP: payrolls,
    unemployment rate, average hourly earnings; PCE: headline and core), with HC3 standard errors.
    122 CPI, 122 NFP and 119 PCE releases (348 with complete \YM\ prices).
  \end{tablenotes}
  \end{threeparttable}
\end{table}

The drift regressions are weaker: with about 120 CPI releases, the coefficient on the headline
surprise alone is not individually significant in any market. The drift is better measured with
the family signal in equation~\eqref{eq:x}, which combines all CPI lines. Its correlation with the
five-minute reaction to CPI is 0.51--0.54 in all four markets, and the 30-minute drift per
standard deviation of the signal is about 5.9 bps in \ES, 8.5 bps in \NQ, 4.4 bps in \YM\ and 0.8 bps in
\ZN. Against a five-minute reaction of 27, 37, 22 and 14 bps per standard deviation, between a fifth
and a quarter of the CPI response in the equity indices remains as a drift after the first minute,
and almost none in \ZN.
Figure~\ref{fig:reaction_drift} shows the same contrast. The trading tests in
Section~\ref{sec:results}, which aggregate many releases with walk-forward signs, are the decisive
evidence on whether this drift is predictable out of sample.

\begin{figure}[htbp]
  \centering
  \includegraphics[width=0.9\linewidth]{nb02b_reaction_vs_drift_4markets}
  \caption{$t$-statistics of the headline surprise (CPI, payrolls, PCE) in regressions of the
  five-minute reaction from the pre-release price (left) and of the drift from one to 30 minutes
  after the release (right), \ES, \NQ, \YM\ and \ZN, 2016--2026. Dashed lines: $\pm 2$.}
  \label{fig:reaction_drift}
\end{figure}

\subsection{Does the pre-release market state matter?}

\citet{Andersen2007} show that announcement effects depend on the state of the economy. We ask a
related question at high frequency: does the same surprise move prices more when markets are
already stressed? We pre-register six state variables measured before the release---20-day
realized volatility, the previous day's intraday volatility, 20-day momentum, the distance from the
60-day high, the overnight return and the previous surprise of the same release---each converted
to a percentile against its trailing one-year distribution, and estimate
\begin{equation}\label{eq:state}
  r_T = a + b\,x_T + c\,s_T + d\,(x_T\times s_T) + \varepsilon_T ,
\end{equation}
where $s_T\in[0,1]$ is the state percentile, so $d$ is the change in the surprise beta from the
calmest to the most stressed state. Each interaction is tested with HC3 and a 2,000-draw
permutation test, corrected for the 48 primary tests with the false-discovery rate of
\citet{Benjamini1995}, and required to have the same sign in 2016--2020 and 2021--2026.

\begin{figure}[htbp]
  \centering
  \includegraphics[width=\linewidth]{fig12_interaction_tstats_CPI_4markets}
  \caption{$t$-statistics of the CPI $\times$ state interaction $d$ in equation~\eqref{eq:state},
  by state variable (rows) and target (columns: reactions and drifts at 1--60 minutes), for \ES,
  \NQ, \YM\ and \ZN.}
  \label{fig:state}
\end{figure}

Figure~\ref{fig:state} summarizes the results. In \ES, CPI surprises hit much harder when recent
volatility is high: the interaction with 20-day volatility has $t$-statistics of 2.9--3.8 on the
reactions at every horizon, and the five-minute beta rises by about 61 bps from the calmest to the
most volatile markets. \NQ\ shows the same pattern slightly weaker ($t$ of 2.2--3.0), and \YM\ almost exactly the same
one: its five-minute CPI beta rises by 58.7 bps from calm to volatile states ($t=3.74$,
permutation $p=0.004$), with the same sign in 2016--2020 and 2021--2026. Three features
limit the economic value of this result. First, the effect is concentrated in the first-minute
jump; on the tradeable drift the $t$-statistics fall to 1.1--2.1. Second, in walk-forward estimation
the interaction was not detectable in real time before 2022: its training $t$-statistic was 0.5 in
2019 and below 2 until 2022. Third, only two \ES\ tests and no \NQ\ test survive the false-discovery
correction; in \YM\ three survive, CPI $\times$ volatility and two NFP $\times$ volatility effects
that come entirely from 2016--2020. On the \YM\ drift the CPI interaction is $+24.6$ bps ($t=2.52$)
but does not survive the correction. In \ZN\ there is no volatility pattern at all. The state variables do predict the size
of the move robustly: in high-volatility states the absolute 30-minute move is about 20 bps larger
in \ES, and 198 (\ES), 154 (\NQ) and 156 (\YM) of 360 magnitude tests survive the correction. We return to this
distinction between direction and size in Section~\ref{sec:mechanisms}.
